% !PRINTED 192
\weberpage{192}
%<<CONT
factors shall be $q$. Then, since each factor contains two
indices,
\weq{19}{\nu + n\mu = 2q.}
Then from (17) and (18), if $P'$ arises from $P$ by the
permutation of $x, \alpha_i$ with $x', \alpha_i'$, it follows that
\weq{20}{{a_0'}^{\mu} P' = a_0^{\mu} r^{n\mu-q} (\gamma x' + \delta)^{\nu} P.}
We now form
\weq{21}{C(x, y, a) = a_0^{\mu} y^{\nu} \Sigma P\left(\frac{x}{y}\right),}
the sum being taken over all the values that one obtains from $P$
when the indices $1, 2, 3 \ldots n$ are permuted with one another
in all possible ways. This sum is then a symmetric function of
the roots of (1), and can therefore be expressed as an integral
rational function of the ratios $a_1 : a_0, a_2 : a_0 \ldots a_n : a_0$,
in such a way that no term exceeds the order $\mu$. The
function $C(x, y, a)$ is therefore an integral homogeneous function of
$a_0, a_1 \ldots a_n$ of order $\mu$, and an integral homogeneous
function of $x, y$ of order $\nu$. Since by (19), (20) it satisfies
the condition
% [sic] The print omits the comma between y' and a' in C(x', y' a').
\weq{22}{C(x', y' a') = r^{\lambda} C(x, y, a),}
\weq{23}{\lambda = n\mu - q, \quad 2\lambda = n\mu - \nu}
it is a \gesperrt{covariant}, and in the particular case
where $\nu = 0$, an \gesperrt{invariant}.

The weight $\lambda$ is, as the second expression shows, always
positive or zero, since $n\mu$ is at least equal to $\nu$. The extreme
value $n\mu = \nu$, or $\lambda = 0$, occurs only when in $P$
each of the indices $1, 2, 3 \ldots n$ occurs joined only with 0,
so that $P$ is a power of $f(x)$ itself.

\begin{center}
\S. 62.

\gesperrt{Binary cubic forms.}
\end{center}

A binary quadratic form gives no invariant
formations other than the discriminant. With these we therefore
do not concern ourselves, and pass at once to the consideration of the
cubic forms
\weq{1}{f(x, y) = a_0 x^3 + a_1 x^2 y + a_2 x y^2 + a_3 y^3}
%CONT>>
